\chapter{Podsumowanie}
\label{cha:Podsumowanie}

\section{Wnioski z realizacji projektu}

W ramach projektu zaimplementowano dwa klasyczne algorytmy sztucznej inteligencji w kontekście praktycznym - jako logikę 
decyzyjną przeciwników w grze komputerowej czasu rzeczywistego. Oba algorytmy działają poprawnie, generują wiarygodne 
zachowanie i stanowią mierzalne wyzwanie dla gracza.
 
\textbf{Co udało się osiągnąć:}
\begin{itemize}
    \item Implementację algorytmu Minimax z przycinaniem alfa-beta na głębokość 4 z zaadaptowaną funkcją oceny 
    pozycji uwzględniającą cztery komponenty (Czebyszew, LOS, power-up, koszt drogi)
    \item Implementację algorytmu A* z heurystyką euklidesową i różnicowymi wagami przejść kafelków
    \item Rozwiązanie kilku nietrywialnych problemów inżynierskich (snap do kafelka po rotacji 34x34, unikanie 
    friendly fire, unik od~sojusznika, optymalizacja THINK\_EVERY)
    \item Bogaty system power-upów wpływający zarówno na gracza, jak i na bota
    \item Dwa odmienne tryby rozgrywki o różnych celach, prezentujące dwie odmienne klasy algorytmów SI
\end{itemize}
 
\textbf{Wnioski algorytmiczne:}
\begin{itemize}
    \item Algorytm Minimax z przycinaniem alfa-beta sprawdza się w grze czasu rzeczywistego nawet przy głębokości 4, 
    pod warunkiem cyklicznego uruchamiania (THINK\_EVERY=15). Pełne przeszukiwanie na każdą klatkę nie jest konieczne ani efektywne.
    \item Wybór heurystyki jest kluczowy. Wartość -2 przy odległości Czebyszewa, +80 za LOS, oraz -60 za symulowany 
    strzał gracza - te konkretne liczby powstały po serii testów, w których obserwowano zachowanie bota i regulowano wagi. 
    Wartości zbyt agresywne sprawiały, że bot nieostrożnie rzucał się na gracza; zbyt defensywne sprawiały, że bot uciekał gdy 
    mógł zaatakować.
    \item A* z heurystyką euklidesową, mimo iż teoretycznie mniej informatywny niż Manhattan dla siatki 4-kierunkowej, działa 
    poprawnie i efektywnie w małej siatce 20x15.
\end{itemize}

\section{Możliwości rozszerzenia projektu}

Projekt w obecnej formie stanowi kompletny prototyp, jednak istnieje szereg kierunków, w których można go rozwinąć:
 
\begin{itemize}
    \item \textbf{Multiplayer sieciowy} - gra dwuosobowa przez sieć, z zachowaniem AI jako trzeciego/czwartego gracza.
    \item \textbf{Reinforcement Learning} - zastąpienie heurystycznej funkcji oceny modelem uczonym. 
    Pozwoliłoby to na AI nauczone "doświadczeniem" grania, niezależne od~ręcznie dobranych wag.
    \item \textbf{Większe i bardziej zróżnicowane mapy} - obecny rozmiar 20x15 kafelków jest stosunkowo mały. 
    Większe mapy uwydatniłyby różnice między algorytmami i potencjalnie wymagałyby głębszego Minimax.
    \item \textbf{Statystyki rozgrywki} - śledzenie wyników, najlepszych czasów, eliminacji - pozwoliłoby na bardziej ilościową 
    ocenę skuteczności AI.
    \item \textbf{Więcej typów wrogów} - różne klasy botów z różnymi parametrami (szybkością, częstością strzału, agresywnością) 
    lub różnymi algorytmami decyzyjnymi.
\end{itemize}